% 精修实验图（pgfplots）——数据来自 docs/BENCH.md / make bench-sweep（CDC 融合后）
\usepackage{pgfplots}
\pgfplotsset{compat=1.18,
  every axis/.append style={
    font=\small,
    label style={font=\small},
    tick label style={font=\footnotesize},
    legend style={font=\footnotesize},
    grid=both,
    grid style={line width=.1pt, draw=gray!25},
    major grid style={line width=.2pt,draw=gray!40},
  }
}

% 图：对数轴 N–延迟（量级对比）
\newcommand{\FigNLatency}{%
\begin{figure}[htbp]
  \centering
  \begin{tikzpicture}
  \begin{axis}[
    width=0.92\textwidth,
    height=0.52\textwidth,
    xlabel={$N$（每 trial 键/文件数）},
    ylabel={中位数延迟（$r\_time$ 相对单位，对数轴）},
    ymode=log,
    xmin=4, xmax=36,
    ymin=1e4, ymax=5e7,
    xtick={8,16,32},
    legend pos=north west,
    legend cell align=left,
  ]
  \addplot+[thick, mark=*, mark size=2.5pt] coordinates {
    (8,33507) (16,38448) (32,53065)
  };
  \addlegendentry{PMFS-KV}
  \addplot+[thick, mark=triangle*, mark size=2.2pt, blue!50!black] coordinates {
    (8,304451) (16,486211) (32,659056)
  };
  \addlegendentry{PMFS-File}
  \addplot+[thick, dashed, mark=square*, mark size=2.2pt, red!70!black] coordinates {
    (8,5684939) (16,11208338) (32,25648605)
  };
  \addlegendentry{xv6 fs+sync}
  \end{axis}
  \end{tikzpicture}
  \caption{随 $N$ 变化的中位数延迟（EMUL，5 trial，CDC 融合后）。对数轴强调相对 xv6 的量级差，不作同等 POSIX 加速比解读。}
  \label{fig:n-latency}
\end{figure}
}

% 图：WA vs N
\newcommand{\FigWaN}{%
\begin{figure}[htbp]
  \centering
  \begin{tikzpicture}
  \begin{axis}[
    width=0.85\textwidth,
    height=0.48\textwidth,
    xlabel={$N$},
    ylabel={$\mathrm{wa\_x100}=\mathrm{phys}\times100/\mathrm{logical}$},
    xmin=4, xmax=36,
    ymin=0, ymax=22000,
    xtick={8,16,32},
  ]
  \addplot+[thick, mark=*, mark size=2.5pt, blue!60!black] coordinates {
    (8,15384) (16,17889) (32,17721)
  };
  \node at (axis cs:20,6000) [align=left, font=\footnotesize] {
    小 value（8\,B）相对 4\,KiB 页\\
    对齐主导放大；融合后随 $N$\\
    近似平台，属页粒度预期
  };
  \end{axis}
  \end{tikzpicture}
  \caption{写放大 $\mathrm{wa\_x100}$ 随 $N$ 变化（EMUL，KV 路径）。高放大属页粒度预期，与 xv6 日志放大不同源。}
  \label{fig:wa-n}
\end{figure}
}

% 图：MAP vs EMUL 柱状（KV 路径）
\newcommand{\FigMapEmul}{%
\begin{figure}[htbp]
  \centering
  \begin{tikzpicture}
  \begin{axis}[
    ybar,
    width=0.75\textwidth,
    height=0.45\textwidth,
    ylabel={KV 中位数延迟（$r\_time$）},
    symbolic x coords={EMUL,MAP},
    xtick=data,
    ymin=0, ymax=140000,
    bar width=18pt,
    nodes near coords,
    every node near coord/.append style={font=\footnotesize},
  ]
  \addplot coordinates {(EMUL,53065) (MAP,113804)};
  \end{axis}
  \end{tikzpicture}
  \caption{$N{=}32$ 时 PMFS-KV 在 EMUL 与 MAP 下的中位数延迟（CDC 融合后）。}
  \label{fig:map-emul}
\end{figure}
}

% 图：文件路径相对吞吐
\newcommand{\FigThr}{%
\begin{figure}[htbp]
  \centering
  \begin{tikzpicture}
  \begin{axis}[
    ybar,
    width=0.85\textwidth,
    height=0.48\textwidth,
    ylabel={相对吞吐（$N\times10^{6}/\mathrm{median}$）},
    symbolic x coords={N8,N16,N32},
    xtick=data,
    ymin=0, ymax=60,
    bar width=12pt,
    legend style={at={(0.98,0.98)},anchor=north east},
  ]
  \addplot coordinates {(N8,26) (N16,32) (N32,48)};
  \addplot coordinates {(N8,1) (N16,1) (N32,1)};
  \legend{PMFS-File,xv6 fs+sync}
  \end{axis}
  \end{tikzpicture}
  \caption{EMUL 下文件 API 相对吞吐（合题「吞吐」指标；非墙钟 MB/s）。}
  \label{fig:thr}
\end{figure}
}
